\section{Summary}\label{summary}

\texttt{unitrootests} provides a unified framework for unit root and
stationarity testing in R, integrating a wide range of classical and
modern tests under a single, consistent interface. The package includes:
quantile ADF tests following Koenker and Xiao (2004) that allow the
order of integration to differ across the conditional distribution;
GARCH-based unit root tests with endogenous structural breaks following
Narayan and Liu (2015) that account for volatility clustering common in
financial and energy price series; and a comprehensive battery of
standard tests (Dickey-Fuller, Phillips-Perron, KPSS, ERS/DF-GLS,
Zivot-Andrews, and Kobayashi-McAleer) guided by the Elder-Kennedy
decision strategy Elder and Kennedy (2001). The package wraps and
extends \texttt{tseries}, \texttt{urca}, and \texttt{strucchange} with
higher-level wrappers that automate model specification and summarize
results.

\section{Statement of Need}\label{statement-of-need}

Determining the order of integration is a fundamental step in time
series econometrics. While R packages such as \texttt{tseries} and
\texttt{urca} provide individual tests, practitioners must manually
coordinate test selection, handle structural break detection, and apply
decision rules across multiple test results. No single R package offers
GARCH-based unit root tests with endogenous breaks,
quantile-distributional unit root tests, and classical tests within a
unified workflow. Researchers in energy economics, finance, and
macroeconomics frequently require all three approaches: classical tests
for baseline inference, GARCH-based tests for volatile series, and
quantile tests to examine tail behavior. \texttt{unitrootests} provides
this integrated toolkit and implements the systematic Elder-Kennedy
decision strategy to guide model selection across test specifications.

\section{Usage}\label{usage}

\subsection{Classical Unit Root Battery with Elder-Kennedy
Strategy}\label{classical-unit-root-battery-with-elder-kennedy-strategy}

\begin{Shaded}
\begin{Highlighting}[]
\FunctionTok{library}\NormalTok{(unitrootests)}

\CommentTok{\# Automated battery: ADF, PP, KPSS, DF{-}GLS, Zivot{-}Andrews}
\NormalTok{result\_ur }\OtherTok{\textless{}{-}} \FunctionTok{urstat}\NormalTok{(y, }\AttributeTok{strategy =} \StringTok{"elder\_kennedy"}\NormalTok{)}
\FunctionTok{summary}\NormalTok{(result\_ur)}
\FunctionTok{print}\NormalTok{(result\_ur)}
\end{Highlighting}
\end{Shaded}

\subsection{Quantile ADF Test}\label{quantile-adf-test}

\begin{Shaded}
\begin{Highlighting}[]
\CommentTok{\# Quantile ADF: test stationarity across quantiles (Koenker \& Xiao, 2004)}
\NormalTok{result\_qadf }\OtherTok{\textless{}{-}} \FunctionTok{qadf}\NormalTok{(y, }\AttributeTok{taus =} \FunctionTok{c}\NormalTok{(}\FloatTok{0.1}\NormalTok{, }\FloatTok{0.25}\NormalTok{, }\FloatTok{0.5}\NormalTok{, }\FloatTok{0.75}\NormalTok{, }\FloatTok{0.9}\NormalTok{),}
                    \AttributeTok{lags =} \DecValTok{2}\NormalTok{)}
\FunctionTok{print}\NormalTok{(result\_qadf)}
\FunctionTok{plot}\NormalTok{(result\_qadf)}
\end{Highlighting}
\end{Shaded}

\subsection{GARCH Unit Root Test with Structural
Breaks}\label{garch-unit-root-test-with-structural-breaks}

\begin{Shaded}
\begin{Highlighting}[]
\CommentTok{\# GARCH{-}based unit root test with endogenous break (Narayan \& Liu, 2015)}
\NormalTok{result\_garch }\OtherTok{\textless{}{-}} \FunctionTok{garchur}\NormalTok{(y, }\AttributeTok{max\_breaks =} \DecValTok{1}\NormalTok{, }\AttributeTok{model =} \StringTok{"GARCH"}\NormalTok{)}
\FunctionTok{summary}\NormalTok{(result\_garch)}
\end{Highlighting}
\end{Shaded}

\section{Implementation}\label{implementation}

\texttt{unitrootests} is built on top of the \texttt{quantreg},
\texttt{tseries}, \texttt{urca}, and \texttt{strucchange} packages. The
quantile ADF test in \texttt{qadf()} follows Koenker and Xiao (2004) by
estimating quantile autoregressions and comparing the quantile-specific
autoregressive coefficient against the null of a unit root, using
critical values from the asymptotic distribution. The \texttt{garchur()}
function implements Narayan and Liu (2015) by fitting a GARCH(1,1) model
with structural break dummies, computing GLS-detrended data, and
applying ADF-type tests on the residuals. The \texttt{urstat()} function
applies the full battery of classical tests and implements the Elder and
Kennedy (2001) decision tree to navigate model specification (trend,
drift, or none) based on sequential test outcomes. All test results are
returned as structured S3 objects with consistent \texttt{print()} and
\texttt{summary()} methods.

\section*{References}\label{references}
\addcontentsline{toc}{section}{References}

\protect\phantomsection\label{refs}
\begin{CSLReferences}{1}{1}
\bibitem[\citeproctext]{ref-Elder2001}
Elder, John, and Peter E. Kennedy. 2001. {``Testing for Unit Roots: What
Should Students Be Taught?''} \emph{Journal of Economic Education} 32
(2): 137--46. \url{https://doi.org/10.1080/00220480109595179}.

\bibitem[\citeproctext]{ref-Koenker2004}
Koenker, Roger, and Zhijie Xiao. 2004. {``Unit Root Quantile
Autoregression Inference.''} \emph{Journal of the American Statistical
Association} 99 (467): 775--87.
\url{https://doi.org/10.1198/016214504000001114}.

\bibitem[\citeproctext]{ref-Narayan2015}
Narayan, Paresh Kumar, and Ruipeng Liu. 2015. {``A Unit Root Model for
Trending Time-Series Energy Variables: An Application to {US} Stock
Prices.''} \emph{Energy Economics} 50: 391--402.
\url{https://doi.org/10.1016/j.eneco.2014.11.021}.

\end{CSLReferences}
